\documentclass[a4paper,12pt]{report}

% Pacotes essenciais
\usepackage[utf8]{inputenc}
\usepackage[T1]{fontenc}
\usepackage[brazil]{babel}
\usepackage{geometry}
\usepackage{titlesec}
\usepackage{hyperref}
\usepackage{enumitem}
\usepackage{xcolor}
\usepackage{fancyhdr}
\usepackage{graphicx}

% Configuração de Margens e Layout
\geometry{
 a4paper,
 total={170mm,257mm},
 left=25mm,
 top=25mm,
}

% Cores e Estilo
\definecolor{primaryColor}{RGB}{50, 0, 90}
\definecolor{secondaryColor}{RGB}{100, 100, 100}

\titleformat{\chapter}{\normalfont\huge\bfseries\color{primaryColor}}{\thechapter}{1em}{}
\titleformat{\section}{\normalfont\Large\bfseries\color{primaryColor}}{\thesection}{1em}{}
\titleformat{\subsection}{\normalfont\large\bfseries\color{secondaryColor}}{\thesubsection}{1em}{}

% Cabeçalho e Rodapé
\pagestyle{fancy}
\fancyhf{}
\rhead{\textbf{PROJETO: SEM TÍTULO} - Game Design Document}
\lhead{v4.0 - Balance \& Lore Update}
\cfoot{\thepage}

\title{\textbf{ENTERTAINMENT GAME PROJECT}\\ \large Game Design Document (GDD)}
\author{Equipe de Desenvolvimento}
\date{\today}

\begin{document}

\maketitle
\tableofcontents

% =========================================================================
% CAPÍTULO 1: VISÃO GERAL
% =========================================================================
\chapter{Visão Geral}

\section{High Concept}
Um \textit{Roguelite} de Ação e Sobrevivência focado na improvisação. O jogador controla um engenheiro civil comum, preso em um planeta vivo que tenta digeri-lo. Sem treinamento militar, ele deve adaptar suas ferramentas de trabalho em armas letais, utilizando a própria biologia do planeta para modificá-las, enquanto busca peças para reparar sua nave e escapar de um ciclo eterno de morte e reconstrução.

\section{A Atmosfera}
O jogo deve transmitir a sensação de \textbf{isolamento claustrofóbico} em um ambiente vasto. Não somos heróis; somos intrusos. O planeta é belo, mas hostil de uma forma biológica: o chão pulsa, as plantas observam e a morte é apenas um processo digestivo interrompido.

% =========================================================================
% CAPÍTULO 2: NARRATIVA E UNIVERSO
% =========================================================================
\chapter{Narrativa e Universo}

\section{O Protagonista: O Engenheiro}
O jogador assume o papel de um funcionário de baixo escalão de uma corporação de logística interplanetária. Seu trabalho consistia em viagens de rotina entre colônias seguras e já catalogadas, realizando manutenção em maquinário pesado.

\begin{figure}[hbt!]
    \centering
    \begin{minipage}[c]{0.48\textwidth}
        \centering
        \includegraphics[width=\linewidth]{AstronautaP.png}
        \caption{Concept Art: Visual Final.}
    \end{minipage}
    \hfill
    \begin{minipage}[c]{0.48\textwidth}
        \centering
        \includegraphics[width=0.65\linewidth]{AstronautaTB.png}
        \caption{Vista Técnica: Frente.}
        \vspace{0.5cm}
        \includegraphics[width=0.65\linewidth]{AstronautaTf.png}
        \caption{Vista Técnica: Costas.}
    \end{minipage}
    \label{fig:protagonista_completo}
\end{figure}

\subsection{Identidade Visual}
\begin{itemize}
    \item \textbf{O Traje (EPI Espacial):} Um macacão pressurizado de cor \textbf{branco gelo} com reforços em \textbf{cinza escuro} nas articulações, luvas e botas. O design prioriza a função sobre a estética, repleto de bolsos pequenos para ferramentas e conectores universais.
    \item \textbf{O Capacete:} O visor possui uma película preta totalmente opaca. Isso serve a dois propósitos: proteger contra a radiação de solda e, narrativamente, permitir que qualquer jogador se projete no personagem (ele não tem rosto definido).
    \item \textbf{Iluminação Diegética:} O traje possui luzes de serviço integradas (ombros e peito) que iluminam os ambientes escuros do jogo naturalmente, dispensando interfaces artificiais de visão noturna.
    \item \textbf{Sem Treinamento de Combate:} Ele nunca segurou uma arma de fogo militar. Todo o seu arsenal é composto por ferramentas de coleta e uma adaga simples de corte.
    \item \textbf{Personalidade:} Prático, curioso e focado em resolver problemas.
\end{itemize}

\section{O Incidente: A Rota Desviada}
A nave de transporte seguia uma rota hiperespacial segura e tediosa. O único companheiro do engenheiro é \textbf{Eptinho}, um Corgi hiperativo com seu próprio traje espacial.
\begin{itemize}
    \item \textbf{O Erro:} Durante um momento de distração, Eptinho perseguiu um reflexo no painel de controle, desengatando o piloto automático e alterando as coordenadas para um setor proibido dos mapas estelares.
    \item \textbf{O Cinturão da Morte:} A nave emergiu do hiperespaço diretamente dentro do campo gravitacional de um planeta não catalogado, cercado por um cinturão de asteroides caótico. Nenhuma nave jamais atravessou esse cinturão intacta. A sobrevivência do Engenheiro à queda foi estatisticamente impossível — um milagre de pura sorte que destruiu completamente \textit{A Nave}, espalhando seus destroços pela superfície.
\end{itemize}

\section{O Planeta: A Entidade Viva}
Este mundo não é apenas um lugar; é um organismo. Tudo nele — a atmosfera, a fauna, a flora e o solo — faz parte de uma única consciência.
\begin{itemize}
    \item \textbf{A Lei da Absorção:} Neste planeta, nada morre verdadeiramente. A matéria orgânica é decomposta, absorvida pelo núcleo e reintegrada em novas formas de vida. É um ciclo perfeito de reciclagem biológica.
    \item \textbf{A Anomalia (O Jogador):} O Engenheiro é feito de uma biologia alienígena a este ecossistema. O planeta tenta absorvê-lo quando ele cai, mas falha. O sistema imunológico planetário não consegue "digerir" a complexidade humana.
    \item \textbf{O Ciclo de Renascimento:} Quando o jogador "morre", o planeta tenta desfazê-lo, mas acaba apenas reconstruindo-o no ponto de entrada (a queda), rejeitando-o como um corpo estranho. O jogador renasce não porque é especial, mas porque é \textbf{indigesto}.
\end{itemize}

\section{A Mecânica do Inventário (A Bolsa)}
Como o planeta absorve tudo o que toca, qualquer equipamento deixado no chão é consumido em segundos.
\begin{itemize}
    \item \textbf{Por que perdemos o Loot:} Ao morrer (ser dissolvido), as armas e modificações externas são comidas pelo planeta.
    \item \textbf{Por que mantemos os Recursos:} A bolsa de coleta do Engenheiro é feita de um polímero sintético inerte hermeticamente fechado. É o único lugar seguro. Tudo o que está \textbf{dentro} da bolsa (minerais, recursos, insetos) sobrevive à reconstrução. O que está \textbf{fora} (armas na mão, armadura) é perdido.
\end{itemize}

\section{Eptinho: O Copiloto e Analista de Dados}
Apesar de sua aparência canina, Eptinho é um ser de inteligência aguçada, atuando como o verdadeiro cérebro por trás dos sistemas de navegação e logística da nave. Ele permanece em segurança na base, mas acompanha o protagonista virtualmente a cada passo.

% ATENÇÃO: coloque o arquivo Eptinho.png na mesma pasta que este .tex
\begin{figure}[h]
    \centering
    \includegraphics[width=0.60\textwidth]{EPTINHO.png}
    \caption{Eptinho: O Copiloto Corgi em seu traje espacial.}
    \label{fig:eptinho}
\end{figure}

\begin{itemize}
    \item \textbf{O Guia da Interface (UI):} Eptinho funciona como o "Oráculo" e sistema operacional da jornada. Ele aparece visualmente como um \textit{pop-up} na lateral da tela (HUD) para notificar o jogador sempre que um novo recurso é catalogado ou quando há itens de coleta pendentes na fase. É ele quem detalha as propriedades dos poderes infundidos e explica os Tiers de raridade dos itens encontrados.
    \item \textbf{Tradutor Universal:} Ele é o responsável por compilar a \textit{Lore} do jogo no banco de dados. Além disso, é o único capaz de compreender a linguagem enigmática do Mercador, servindo como intérprete durante as trocas comerciais — uma habilidade que se aprimora a cada encontro (ver Capítulo~\ref{cap:mercador}, Seção~\ref{sec:traducao}).
    \item \textbf{Comunicação:} Possui um link de áudio direto com o capacete do Engenheiro. Embora se expresse através de latidos, o sistema de tradução do traje converte instantaneamente o áudio em texto no visor do jogador, permitindo que Eptinho explique a história do planeta e forneça suporte tático em tempo real.
\end{itemize}

% =========================================================================
% CAPÍTULO 3: SISTEMA DE ATRIBUTOS E ESTATÍSTICAS
% =========================================================================
\chapter{Sistema de Atributos e Estatísticas}

Para permitir uma profundidade estratégica nas \textit{builds}, o jogo utiliza um conjunto extenso de variáveis que podem ser modificadas por Infusões (Partes do Corpo), Itens do Mercador e Upgrades da Base.

\section{Atributos Ofensivos}
\begin{description}
    \item[Attack Speed (Melee):] Reduz o intervalo entre os golpes de armas de curto alcance.
    \item[Crit Chance:] A probabilidade percentual de um ataque causar dano amplificado.
    \item[Crit Multiplier:] O fator de multiplicação do dano quando um acerto crítico ocorre (Ex: 0.15X, 0.25X).
    \item[Weapon Range (Melee):] Aumenta o tamanho da \textit{Hitbox} ou a extensão física do golpe para armas de corpo-a-corpo.
    \item[Weapon Range (Projectiles):] Aumenta a distância máxima que um projétil viaja antes de desaparecer ou perder eficácia.
    \item[Knockback:] A força física aplicada ao inimigo ao ser atingido, empurrando-o para trás. Essencial para controle de multidão.
    \item[Piercing:] O número de inimigos que um projétil ou golpe pode atravessar antes de ser destruído.
    \item[Bounce (Chance/Qtd):] Para armas de longo alcance, define a chance ou quantidade de vezes que o projétil ricocheteia para um inimigo próximo após o impacto.
    \item[MultiShot (Chance):] A probabilidade de disparar projéteis adicionais simultaneamente sem consumir recursos extras.
    \item[Spread:] Define o ângulo de dispersão dos disparos. Valores menores indicam maior precisão; valores maiores indicam uma área de cobertura mais ampla (estilo escopeta).
\end{description}

\section{Atributos Defensivos}
\begin{description}
    \item[Armor Regen:] A velocidade com que a barra de Armadura (Escudo) se regenera após o tempo de recarga.
    \item[Dodge:] A chance percentual de anular completamente o dano de um ataque recebido.
    \item[Damage Negation:] Uma redução fixa ou percentual em todo dano recebido (Mitigação direta).
    \item[Thorns:] Dano devolvido automaticamente ao atacante quando o jogador é atingido por contato físico.
\end{description}

\section{Atributos de Mobilidade}
\begin{description}
    \item[Dash Cooldown:] Reduz o tempo de espera necessário para utilizar a esquiva novamente.
    \item[Dash Counts:] Aumenta o número máximo de cargas de Dash que o jogador pode armazenar.
    \item[Dash Invulnerability:] Aumenta a duração da janela de invencibilidade (\textit{i-frames}) durante a animação da esquiva.
    \item[Speed:] Aumenta a velocidade de movimento base do personagem.
\end{description}

\section{Status e Economia (Mercador/Coleta)}
\begin{description}
    \item[Luck (Essence):] Aumenta a quantidade de \textit{Essência da Vida} obtida por orbe coletado.
    \item[Luck (Drops):] Aumenta a probabilidade global de inimigos soltarem itens (\textit{Partes do Corpo}) e a chance de virem com Tiers superiores.
    \item[Magnet:] Aumenta o raio de coleta automática para atrair Essência e itens distantes.
\end{description}

\section{Atributos da Habilidade Suprema (ULT)}
\begin{description}
    \item[Cooldown Reduction (ULT):] Reduz o tempo de recarga da Habilidade Suprema.
    \item[Ult Charge with Kills:] Permite que a eliminação de inimigos reduza o cooldown ativo da ULT (Recarga por abate).
    \item[Ult Radius:] Aumenta a área de efeito (AoE) de explosões ou buffs da Suprema.
    \item[Ult Damage:] Um multiplicador de dano exclusivo para a habilidade especial.
    \item[Ult Effect Duration:] Aumenta o tempo de permanência de efeitos criados pela ULT (ex: duração de uma poça de ácido ou de uma tempestade elétrica).
    \item[Ult Buff Potency:] Para Supremas que concedem \textit{buffs} (ex: Sobrecarga do Rifle), aumenta os status concedidos durante a ativação.
\end{description}

\section{Core Loop: O Ciclo de Infusão}
O núcleo da jogabilidade abandona o sistema tradicional de níveis (XP) por um sistema de \textbf{Economia e Crafting em Tempo Real}. O ciclo se define por quatro etapas constantes:

\begin{enumerate}
    \item \textbf{Combate:} O jogador enfrenta inimigos usando um dos \textbf{6 Arquétipos de Armas} (ver Capítulo~\ref{cap:arsenal}), que definem seu estilo de jogo.
    \item \textbf{Coleta (Loot \& Currency):} Inimigos derrotados sempre soltam \textit{Essência da Vida} (moeda) e têm uma chance probabilística de soltar \textit{Partes do Corpo} (matéria-prima).
    \item \textbf{Decisão Estratégica:} Ao pegar uma Parte do Corpo, o jogador decide instantaneamente:
    \begin{itemize}
        \item \textbf{Infundir:} Gastar Essência para acoplar a parte à arma, ganhando poder imediato.
        \item \textbf{Descartar/Ignorar:} Economizar Essência para um drop futuro de Tier maior.
    \end{itemize}
    \item \textbf{Adaptação:} A arma do jogador evolui organicamente durante a run, tornando-se uma quimera de tecnologias humanas e biologia alienígena.
\end{enumerate}

\section{Ecossistema de Recursos}

\subsection{1. Essência da Vida (A Moeda Universal)}
\begin{itemize}
    \item \textbf{Origem:} Drop garantido de 100\% dos inimigos. A quantidade escala conforme a força do inimigo.
    \item \textbf{Função:} Usada para realizar infusões na arma durante a run e comprar itens no Mercador.
    \item \textbf{Lore:} A energia vital pura do planeta, capaz de moldar a matéria e reescrever o código genético dos materiais.
\end{itemize}

\subsection{2. Partes do Corpo (O Loot)}
\begin{itemize}
    \item \textbf{Origem:} Drop probabilístico (RNG) dos inimigos.
    \item \textbf{Natureza ``Não Identificada'':} No chão e no inventário, a parte exibe apenas seu \textbf{Nome} (ex: "Mandíbula de Inseto") e seu \textbf{Tier} (1 a 4). O jogador \textit{não sabe} o efeito exato até gastar essência para infundi-la.
\end{itemize}

\subsubsection{Tiers de Raridade}
\begin{description}
    \item[Tier 1 (Comum)] \hfill \\
    Melhorias numéricas baixas (ex: +3\% Dano, +2\% Velocidade de Ataque). Amplificação máxima de +5\%.
    \item[Tier 2 (Incomum)] \hfill \\
    Melhorias numéricas robustas (ex: +8\% Dano, +5\% Chance Crítica). Amplificação máxima de +8\%.
    \item[Tier 3 (Raro)] \hfill \\
    Melhorias numéricas potentes (ex: +10\% Dano, +5\% Velocidade de Ataque, +8\% Chance Crítica). Amplificação máxima de +10\%.
    \item[Tier 4 --- Lendário] \hfill \\
    Alterações drásticas de gameplay (ex: ``Sua arma não usa mais munição, mas consome vida para atirar''). \textit{Garantido por Guardiões.} Muito caros para infundir; amplificação possível com custo exacerbado.
\end{description}

\section{Progressão durante a Run (Sistema de Infusão)}

\subsection{Mecânica de Infusão}
O jogador pode carregar um número limitado de Partes do Corpo. A qualquer momento fora de combate, ele pode abrir o menu de Infusão:
\begin{enumerate}
    \item Seleciona uma Parte do Corpo coletada.
    \item Paga um custo em \textbf{Essência da Vida}.
    \item A Parte é consumida e um atributo aleatório (baseado no Tier e tipo da parte) é acoplado à arma permanentemente naquela run.
\end{enumerate}

\subsection{Inflação de Custo}
Para impedir que o jogador fique infinitamente forte farmando inimigos fracos, o sistema aplica uma \textbf{inflação progressiva}: cada infusão realizada aumenta o custo da próxima. O peso acumulado de todas as infusões ($P_{total}$) determina o multiplicador de custo.

A especificação matemática completa — fórmula, pesos por Tier e exemplos de escalonamento — está documentada no Capítulo~\ref{cap:economia}, Seção~\ref{sec:inflacao}.

\subsection{Amplificação (Overcharge)}
Ao realizar uma infusão e revelar o atributo, o jogador tem uma breve janela para ``Sobrecarregar'' a infusão, gastando uma quantia extra e massiva de Essência para aprimorar o atributo recém-adquirido.

\subsection{Exemplos de Tipos de Melhorias}
\begin{itemize}
    \item \textbf{Numéricas:} Aumento de Dano, Velocidade de Ataque, Alcance, Chance de Crítico.
    \item \textbf{Mecânicas (Tier 4):}
    \begin{itemize}
        \item \textit{Fio Infinito:} Cada acerto consecutivo aumenta o dano em 0.1\% (reseta ao tomar dano).
        \item \textit{Toxina Viral:} Acertos aplicam veneno que se espalha para inimigos próximos se o alvo morrer infectado.
        \item \textit{Vampirismo de Essência:} Chance de dropar Essência extra ao matar inimigos com Crítico.
    \end{itemize}
\end{itemize}

% ATENÇÃO: substitua 'grafico_matlab.png' pelo nome real do arquivo exportado do MATLAB
\begin{figure}[h]
    \centering
    \includegraphics[width=1\textwidth]{grafico.png}
    \caption{Curva de Custo de Infusão por Composição de Build (gerado no MATLAB).}
\end{figure}

% =========================================================================
% CAPÍTULO 4: O MERCADOR
% =========================================================================
\chapter{O Mercador: A Válvula de Escape}\label{cap:mercador}

\begin{figure}[h]
    \centering
    \includegraphics[width=0.5\textwidth]{Mercador (1).png}
    \caption{O Mercador: A figura misteriosa que negocia com sangue.}
\end{figure}

\begin{figure}[h]
    \centering
    \begin{minipage}{0.4\textwidth}
        \centering
        \includegraphics[width=\linewidth]{Mercador (2).png}
        \caption{Detalhe Lateral: Frascos de Essência.}
    \end{minipage}
    \hfill
    \begin{minipage}{0.4\textwidth}
        \centering
        \includegraphics[width=\linewidth]{Mercador (3).png}
        \caption{Vista Traseira: Talismãs e manto.}
    \end{minipage}
\end{figure}

Ocasionalmente, escondido em becos sem saída entre os biomas, o jogador \textbf{poderá} encontrar uma figura misteriosa conhecida apenas como ``O Mercador''. O encontro \textbf{não é garantido}: existe uma chance probabilística de sua presença a cada \textit{run}, e algumas jornadas podem transcorrer inteiramente sem que ele apareça (ver Seção~\ref{sec:regras_geracao}).

\section{Lore e Aparência}
\begin{itemize}
    \item \textbf{Visual:} Ele veste um manto roxo desgastado com detalhes em couro e um chapéu de aba larga e pontudo que oculta completamente seu rosto, deixando visíveis apenas olhos que brilham na escuridão.
    \item \textbf{Comunicação:} Ele fala uma língua antiga e incompreensível para humanos. Toda a negociação é intermediada por \textbf{Eptinho}, que traduz os sussurros do Mercador através do HUD do jogador — mas essa tradução não é perfeita desde o início.
    \item \textbf{Função no Mundo:} Ele atua como a ``válvula de escape'' do sistema econômico do jogo, permitindo que o jogador corrija erros de \textit{build} ou assuma riscos mortais em troca de poder absoluto.
\end{itemize}

\section{A Tradução Progressiva}\label{sec:traducao}
A língua do Mercador é incompreensível para humanos, mas não para Eptinho — desde que ele tenha tempo suficiente para analisá-la. A cada encontro durante uma mesma \textit{run}, Eptinho acumula mais dados linguísticos e consegue traduzir uma porção maior do discurso.

\begin{description}
    \item[1º Encontro:] Eptinho capta apenas fragmentos e intenções básicas. A maior parte do diálogo aparece como texto corrompido ou símbolos indecifrávels no HUD. O jogador entende o suficiente para identificar os quatro serviços, mas o tom e o contexto do Mercador ficam ocultos.

    \item[2º Encontro:] Com mais dados compilados, Eptinho consegue traduzir aproximadamente metade do discurso. Nuances começam a emergir: o Mercador pode revelar detalhes sobre os pactos ou tecer comentários sobre o estado atual do jogador.

    \item[3º Encontro em diante:] A tradução está quase completa. O Mercador começa a se revelar, compartilhando fragmentos de sua própria perspectiva sobre o planeta, o ciclo de renascimento e a natureza do Engenheiro. Esta é a camada de \textit{Lore} mais profunda do jogo, acessível somente a \textit{runs} que encontram o Mercador múltiplas vezes.
\end{description}

\textit{Nota de Design:} Este sistema incentiva o jogador a buscar o Mercador ativamente em cada run. Encontrá-lo uma única vez é vantajoso economicamente; encontrá-lo várias vezes desbloqueia a narrativa oculta.

\section{Regra de Encontro: Um Serviço por Visita}
\begin{quote}
    \textbf{Regra Fundamental:} Em cada encontro com o Mercador, o jogador pode utilizar \textbf{apenas um} dos quatro serviços disponíveis. Após a transação, o Mercador encerra a negociação e desaparece até o próximo encontro.
\end{quote}

Esta restrição impõe uma decisão estratégica de alto custo de oportunidade: escolher o \textit{Câmbio de Sangue} significa abrir mão da \textit{Cirurgia de Remoção} naquele momento, e vice-versa.

\section{Os Quatro Pilares de Troca}

\subsection{1. Câmbio de Sangue (Recurso de Desespero)}
\begin{description}
    \item[A Troca:] O jogador sacrifica permanentemente uma porcentagem de sua \textbf{Vida Máxima} (irreparável durante a run).
    \item[O Ganho:] Recebe uma injeção massiva e imediata de \textbf{Essência da Vida}.
    \item[Função Estratégica:] Permite comprar um item Lendário ou realizar uma infusão crítica de última hora, salvando uma run mal gerenciada economicamente.
\end{description}

\subsection{2. Cirurgia de Remoção (Otimização)}
\begin{description}
    \item[A Troca:] Um custo moderado de Essência.
    \item[O Ganho:] O Mercador extrai cirurgicamente uma ``Parte do Corpo'' já acoplada à arma, removendo permanentemente seu efeito.
    \item[Recálculo de Inflação:] Ao remover uma parte, o \textbf{peso de inflação} ($P$) daquele item é imediatamente \textbf{subtraído} do $P_{total}$ acumulado. Isso reduz o custo das próximas infusões proporcionalmente ao Tier do item removido. O Mercador efetivamente ``volta o relógio'' da economia. Ver Seção~\ref{sec:inflacao} (Capítulo~\ref{cap:economia}) para os valores de $P$ por Tier.
    \item[Nota:] O item extraído é destruído — ele não retorna ao inventário do jogador.
\end{description}

\subsection{3. Artefatos Refinados \& Únicos}
\begin{description}
    \item[A Troca:] Um custo altíssimo de Essência (muito superior a uma infusão comum).
    \item[O Ganho:] \textbf{Partes Identificadas} — itens de Tier alto onde o jogador \textit{vê} exatamente o atributo antes de comprar.
    \item[Regra de Inflação:] Itens comprados aqui vão para o inventário do jogador. O $P_{total}$ \textbf{não é afetado no momento da compra}, apenas no momento em que o jogador optar por infundir o item normalmente. Isso significa que a compra em si não encarece as próximas infusões.
\end{description}

\subsection{4. As Maldições (O Tarô Proibido)}
\begin{description}
    \item[A Troca:] Custa sempre \textbf{Vida Máxima}.
    \item[O Ganho:] Um poder passivo avassalador que dura até o fim da run, sempre atrelado a uma penalidade severa. Transforma o estilo de jogo imediatamente.
\end{description}

\subsubsection{Exemplos de Cartas de Maldição}
\begin{itemize}
    \item \textbf{A Ganância:} Inimigos têm chance dobrada de dropar loot e Essência. \textit{Maldição:} Inimigos causam 50\% mais dano.
    \item \textbf{O Frenesi:} Velocidade de Ataque e Movimento aumentam em 40\%. \textit{Maldição:} Armadura reduzida a 0 permanentemente.
    \item \textbf{O Parasita:} Vampirismo (regenera vida ao matar). \textit{Maldição:} Necrose acelerada — a vida decai se o jogador ficar mais de 5 segundos sem matar. Força estilo agressivo e sem pausas.
\end{itemize}

% =========================================================================
% CAPÍTULO 5: ARSENAL (FERRAMENTAS ADAPTADAS)
% =========================================================================
\chapter{Arsenal (Ferramentas Adaptadas)}\label{cap:arsenal}

O jogador não possui armas militares. Ele luta com o que tem na maleta de ferramentas. A \textbf{``Classe''} do jogador é definida pela \textbf{arma escolhida no início da run}. Existem 6 arquétipos organizados por Alcance e Cadência. Cada um possui uma \textbf{ULT} (Ataque Especial com Cooldown) que define o momento de pico do estilo de combate.

\section{Curto Alcance / Alta Cadência}
\begin{description}
    \item[Cortador de Plasma Portátil (Adaga):] Originalmente usado para cortar chapas de metal. Alta cadência, baixo alcance.
    \begin{itemize}
        \item \textit{Estilo:} \textit{Hit \& Run}. Agressividade e alto DPS por aproximação.
        \item \textit{ULT:} (?).
    \end{itemize}
\end{description}

\section{Curto Alcance / Baixa Cadência}
\begin{description}
    \item[(?) :] (?).
    \begin{itemize}
        \item \textit{Estilo:} (?).
        \item \textit{ULT:} (?).
    \end{itemize}
\end{description}

\section{Médio Alcance / Alta Cadência}
\begin{description}
    \item[(?) :] (?).
    \begin{itemize}
        \item \textit{Estilo:} (?).
        \item \textit{ULT:} (?).
    \end{itemize}
\end{description}

\section{Médio Alcance / Baixa Cadência}
\begin{description}
    \item[(?) :] (?).
    \begin{itemize}
        \item \textit{Estilo:} (?).
        \item \textit{ULT:} (?).
    \end{itemize}
\end{description}

\section{Longo Alcance / Alta Cadência}
\begin{description}
    \item[(?) :] (?).
    \begin{itemize}
        \item \textit{Estilo:} (?).
        \item \textit{ULT:} (?).
    \end{itemize}
\end{description}

\section{Longo Alcance / Baixa Cadência}
\begin{description}
    \item[(?) :] (?).
    \begin{itemize}
        \item \textit{Estilo:} (?).
        \item \textit{ULT:} (?).
    \end{itemize}
\end{description}

% =========================================================================
% CAPÍTULO 6: BESTIÁRIO DE INIMIGOS
% =========================================================================
\chapter{Bestiário de Inimigos}\label{cap:bestiario}

As criaturas do bioma cristalino são manifestações físicas do sistema imunológico do planeta. Cada classe define o custo de spawn e as restrições táticas dentro das ondas de combate. Os inimigos concretos listados neste capítulo estão implementados e ativos no projeto.

% =========================================================================
\section{Tabela de Classes de Spawn (Sistema de Budget)}

\begin{center}
\begin{tabular}{|l|c|l|}
    \hline
    \textbf{Classe} & \textbf{Custo} & \textbf{Restrições} \\
    \hline
    Mob Menor            & 1 ponto   & Máx. 50--70\% dos pontos da onda \\
    Atirador             & 4 pontos  & Sempre em par (2 por 4 pts) \\
    Tanque               & 4 pontos  & Máx. 4 por sala (limite global) \\
    Suporte (Sintonizador) & 5 pontos  & Máx. 2 por sala \\
    Estrutura (Totem)    & 6 pontos  & Máx. 1 por sala; crânios não consomem budget \\
    Elite                & 10 pontos & Máx. 1 por sala (limite global) \\
    \hline
\end{tabular}
\end{center}

% =========================================================================
\section{Inimigos Implementados --- Bioma 1 (Cristalizado)}

% -------------------------------------------------------------------------
\subsection{Aranha (\textit{Spider\_AI}) --- Mob Menor}
% ATENÇÃO: adicione o arquivo de imagem na pasta do .tex
\begin{figure}[h]
    \centering
    \includegraphics[width=0.45\textwidth]{Inimigo_Aranha.png}
    \caption{Aranha: Mob Menor ágil com leap attack e recuo tático.}
\end{figure}

\begin{description}
    \item[Lore:] Fragmento vivo do tecido do planeta. Ágil e predatória, é o primeiro organismo a responder quando a anomalia humana invade o território.
    \item[Comportamento:]
    \begin{itemize}
        \item Aproxima-se do jogador até o range de salto (\texttt{leapRange = 5m}).
        \item \textbf{Passo Procedural:} Movimenta-se usando um ciclo de pernas procedural 100\% dinâmico, suportando padrões de caminhada alternados (\textit{Tripod Gait} de 6 pernas ou \textit{Wave Gait} sequencial).
        \item \textbf{Leap Attack:} Salta em trajetória livre com força de impulso (\texttt{leapForce = 12}); causa dano por \textit{OverlapSphere} ao aterrissar (\texttt{leapDamage = 10}). Durante o salto, suas garras/mandíbulas se abrem e fecham proceduralmente.
        \item \textbf{Recuo Tático:} 30\% de chance de saltar para trás após um hit bem-sucedido, dificultando a punição do jogador.
        \item \textbf{VFX:} Rastro de fantasmas visual (\textit{after-images}) via \texttt{SpiderDashVFX} durante o pulo de ataque.
    \end{itemize}
    \item[Stats:] Velocidade \texttt{6} --- Dano \texttt{10} --- Cooldown do salto \texttt{2s} --- Cooldown de recuo \texttt{5s}
\end{description}

% -------------------------------------------------------------------------
\subsection{Goblin (\textit{GoblinAI\_Transform}) --- Atirador / Melee Situacional}\label{sec:goblin}
\begin{figure}[h]
    \centering
    \includegraphics[width=0.35\textwidth]{Inimigo_Goblin.png}
    \caption{Goblin: Mineiro subterrâneo que arremessa bombas à distância e ataca com picareta quando encurralado.}
\end{figure}

\begin{description}
    \item[Lore:] Os Goblins são uma raça de seres subterrâneos que habitam as profundezas do planeta, minerando cristais com propriedades explosivas. Por passarem toda a vida nas cavernas sem exposição à claridade, seus olhos atrofiaram --- são essencialmente cegos ou hipersensíveis à luz, cobrindo sempre o rosto com um pano grosseiro. Fisicamente são atarracados e esguios: a combinação de baixa gordura e musculatura densa de anos de mineração os torna visualmente frágeis, mas surpreendentemente fortes. Primitivos em cultura, carregam apenas uma bolsa de couro com cristais e a picareta que os define.
    \item[Comportamento:]
    \begin{itemize}
        \item Mantém distância ideal de \texttt{12m} realizando strafe lateral contínuo.
        \item \textbf{Bomba Parabólica (Ranged):} Arremessa \texttt{BombaExplosiva} a cada \texttt{2.8s}. Pisca em vermelho e detona em \texttt{3s} (\texttt{dano = 40}, \texttt{raio = 2m}, knockback).
        \item \textbf{Fuga:} Se o jogador entra em \texttt{6m}, tenta fugir a \texttt{velocidade = 11}.
        \item \textbf{Ataque de Picareta (Melee):} Se o jogador entra em \texttt{2.5m} --- distância menor que a de fuga --- o Goblin \textbf{interrompe a fuga} e desfere um golpe de picareta. Para completamente, levanta a ferramenta (wind-up \texttt{0.35s}) e causa \texttt{20} de dano em raio de \texttt{1.8m} à sua frente. Cooldown de \texttt{1.5s} entre golpes.
        \item Aparece em par no sistema de spawn (2 unidades por 4 pontos).
    \end{itemize}
    \item[Prioridade de estados:] \texttt{MeleePickaxe} ($\leq$2.5m) $>$ \texttt{Fugir} ($<$6m) $>$ \texttt{Atacar Ranged} ($\leq$12m) $>$ \texttt{Perseguir}
    \item[Stats:] Bomba: Dano \texttt{40} --- Raio \texttt{2m} (buffado \texttt{4m}) --- Picareta: Dano \texttt{20} --- Cooldown \texttt{1.5s}
    \item[Buff do Sintonizador:] Velocidade +15\%, intervalo de bomba $\div$2, raio da bomba dobra.
\end{description}

% -------------------------------------------------------------------------
\subsection{Golem (\textit{Golem\_AI}) --- Tanque}
\begin{figure}[h]
    \centering
    \includegraphics[width=0.45\textwidth]{Inimigo_Golem.png}
    \caption{Golem: Tanque com melee pesado e stun em área após combo.}
\end{figure}

\begin{description}
    \item[Lore:] A consolidação máxima da estrutura mineral do bioma. Lento demais para ser ignorado, poderoso o suficiente para terminar qualquer descuido.
    \item[Comportamento:]
    \begin{itemize}
        \item \textbf{Animação Procedural:} Move-se lentamente em direção ao jogador (\texttt{moveSpeed = 2}), com oscilações pesadas, inércia de corpo, spine corcunda e inclinações de quadril simuladas proceduralmente.
        \item \textbf{Melee (Ataque Rotatório):} A \texttt{3m} de distância, golpeia com wind-up telegrafado. Os braços inclinam-se para dentro e giram em rotações simétricas opostas.
        \begin{itemize}
            \item \textbf{Convergência das Mãos:} Os braços convergem inclinando-se para baixo ($30^\circ$) e para dentro ($20^\circ$) nas extremidades das mãos, garantindo que o golem atinja o jogador (que é mais baixo e menor).
            \item \textbf{Limite de Rotação:} Sem buff, o ataque melee realiza apenas meio giro ($180^\circ$ de varredura). Sob o efeito do buff do Sintonizador, os braços executam 2 giros completos ($720^\circ$ de rotação contínua).
        \end{itemize}
        \item \textbf{Combo Stun:} Após \textbf{2 acertos de melee}, ativa um stun em área (\texttt{raio = 4m}) na posição do jogador.
        \item \textbf{Stun Proativo:} Se o cooldown do melee estiver ativo e o jogador estiver em \texttt{10m}, usa o stun como ataque primário.
        \item \textbf{VFX do Stun:} Efeito URP-emissivo moderno que projeta anéis de impacto expansivos e translúcidos, 4 rachaduras angulares no terreno que crescem para fora, um pilar de energia central com contração de largura, e partículas em três camadas (faíscas velozes, pedregulhos de cristal com gravidade, e poeira expansiva rente ao chão).
    \end{itemize}
    \item[Stats:] Dano melee \texttt{35} --- Cooldown melee \texttt{3s} --- Stun \texttt{1.5s} --- Cooldown stun \texttt{12s}
    \item[Buff do Sintonizador:] Velocidade +30\%, cooldown do stun $\times$0.7, e ataque giratório com 2 voltas completas ($720^\circ$). \textit{Nota: Ao destruir o Sintonizador que o acompanha, o Golem é instantaneamente desbuffado, revertendo imediatamente sua velocidade, cooldowns, cor da barra de vida (de ciano para vermelho) e limitando seu ataque giratório de volta para os $180^\circ$ padrão.}
\end{description}

% -------------------------------------------------------------------------
\subsection{Pedra Mágica (\textit{MagicStone\_AI}) --- Elite Ranged}
\begin{figure}[h]
    \centering
    \includegraphics[width=0.45\textwidth]{Inimigo_MagicStone.png}
    \caption{Pedra Mágica: Elite flutuante com raio vertical telegrafado e teleporte de emergência.}
\end{figure}

\begin{description}
    \item[Lore:] Fragmento de consciência planetária condensado. Flutua sobre o campo de batalha e golpeia o intruso com energia pura do núcleo.
    \item[Comportamento:]
    \begin{itemize}
        \item Mantém órbita entre \texttt{12--15m} do jogador, circulando continuamente com flutuação sinusoidal vertical.
        \item \textbf{Skybeam Attack:} A cada \texttt{5s}, telegrafeia com \texttt{AttackMarkerEffect} por \texttt{2.5s} na posição atual do jogador, depois dispara um raio vertical (\texttt{AttackBeam}).
        \item \textbf{Teleporte de Emergência:} Se o jogador se aproximar a menos de \texttt{4m}, teleporta para \texttt{15--20m}, respeitando os limites da sala (\texttt{roomBounds}).
    \end{itemize}
    \item[Stats:] Intervalo de ataque \texttt{5s} --- Telegrafagem \texttt{2.5s} --- Cooldown teleporte \texttt{30s}
    \item[Buff do Sintonizador:] Intervalo de ataque $\div$2, velocidade de órbita +50\%.
\end{description}

% -------------------------------------------------------------------------
\subsection{Enxame de Cristais (\textit{ShardSwarm\_AI}) --- Elite}
\begin{figure}[h]
    \centering
    \includegraphics[width=0.45\textwidth]{Inimigo_ShardSwarm.png}
    \caption{Enxame de Cristais: Ataca em grupo, se divide ao tomar burst de dano e explode ao morrer.}
\end{figure}

\begin{description}
    \item[Lore:] O planeta em sua forma mais primitiva — fragmentos orbitais animados por impulso de destruição, sem forma definitiva e impossíveis de conter.
    \item[Comportamento:]
    \begin{itemize}
        \item Fragmentos orbitam o núcleo; conforme o HP cai, fragmentos somem visualmente.
        \item \textbf{Swarm Attack:} Dentro de \texttt{5m}, todos os fragmentos voam simultaneamente ao jogador (\texttt{dano = 20 combinado}, cooldown \texttt{2s}).
        \item \textbf{Divisão (Split):} Se receber $>$30\% do HP máximo de uma vez, clona-se ao lado com metade do HP. O clone não pode se dividir.
        \item \textbf{Explosão de Morte:} Ao morrer, explode em \texttt{raio = 3m} causando \texttt{dano = 15}.
    \end{itemize}
    \item[Stats:] Dano combinado \texttt{20} --- Dano individual \texttt{5}/fragmento --- Gatilho split \texttt{30\%} do HP em burst
    \item[Buff do Sintonizador:] Dano +50\%.
\end{description}

% -------------------------------------------------------------------------
\subsection{Vigia de Cristal (\textit{CrystalWatcher\_AI}) --- Elite Estático}
\begin{figure}[h]
    \centering
    \includegraphics[width=0.45\textwidth]{Inimigo_CrystalWatcher.png}
    \caption{Vigia de Cristal: Imóvel. Laser giratório que inverte conforme o jogador dasha.}
\end{figure}

\begin{description}
    \item[Lore:] Sentinela do território. Não persegue — não precisa. O laser emitido transforma qualquer rota de fuga em cálculo de risco.
    \item[Comportamento:]
    \begin{itemize}
        \item Estático. Ativa quando o jogador entra em \texttt{20m}.
        \item \textbf{Ciclo:} Carregamento (\texttt{0.5s}) $\rightarrow$ Laser ativo (\texttt{4s}) $\rightarrow$ Cooldown (\texttt{2s}).
        \item \textbf{Laser Giratório:} Gira pelo caminho mais curto até o jogador (\texttt{30°/s}). Se o jogador dasha para o lado oposto, o laser \textbf{inverte a rotação automaticamente}.
        \item \textbf{Dano por Tick:} \texttt{15} de dano a cada \texttt{0.3s} de contato (\textit{SphereCastAll}, \texttt{laserWidth = 0.5m}).
        \item \textbf{Visual:} Dois \texttt{LineRenderers} — núcleo fino branco-roxo + glow externo roxo pulsante.
    \end{itemize}
    \item[Stats:] Dano/tick \texttt{15} --- Intervalo \texttt{0.3s} --- Rotação \texttt{30°/s} --- Alcance \texttt{20m}
    \item[Buff do Sintonizador:] Carregamento $\div$2, rotação do laser +50\% (\texttt{45°/s}).
\end{description}

% =========================================================================
\section{Inimigos de Suporte e Controle}

% -------------------------------------------------------------------------
\subsection{Sintonizador de Cristal (\textit{CrystalTuner})}
\begin{figure}[h]
    \centering
    \includegraphics[width=0.45\textwidth]{Inimigo_CrystalTuner.png}
    \caption{Sintonizador: Voa e amplifica aliados via feixes de energia. Eliminá-lo é prioridade.}
\end{figure}

\begin{description}
    \item[Lore:] O sistema nervoso do planeta em forma ambulante. Liga-se a outros organismos via feixes de energia pura, elevando sua capacidade combativa.
    \item[Comportamento:]
    \begin{itemize}
        \item Voa a \texttt{1.5m} de altura e posiciona-se próximo aos aliados.
        \item Foge do jogador se ele entrar em \texttt{8m}.
        \item Conecta-se a até \textbf{3 inimigos simultâneos} dentro de \texttt{20m} com feixes visuais ondulados (Perlin noise), concedendo buffs individualizados por tipo.
        \item \textbf{Feedback de Conexão:}
        \begin{itemize}
            \item \textbf{Point Light:} Possui uma luz que se ativa apenas quando conectado a aliados, variando e pulsando sua intensidade continuamente. Desativa-se quando não há alvos conectados.
            \item \textbf{Rotação do Modelo:} O corpo de cristal gira continuamente. Em estado ocioso (\textit{idle}), rotaciona lentamente a $30^\circ/\text{s}$. Quando ativado e transmitindo buffs, acelera a rotação para $180^\circ/\text{s}$.
        \end{itemize}
        \item \textbf{Morte/Desativação Limpa:} Ao morrer ou ser desativado, o Sintonizador remove imediatamente os buffs de todos os seus alvos, destrói os feixes de conexão, desativa sua Point Light e interrompe a aceleração de rotação.
    \end{itemize}
    \item[Stats:] Máx. alvos \texttt{3} --- Range \texttt{20m} --- Velocidade \texttt{4.5} --- Fuga abaixo de \texttt{8m}
\end{description}

% -------------------------------------------------------------------------
\subsection{Totem (\textit{TotemSpawner}) --- Estrutura}
\begin{figure}[h]
    \centering
    
\includegraphics[width=0.45\textwidth]{Inimigo_Totem.jpg}
    \caption{Totem: Estrutura fixa que reanima crânios de Goblins mortos. Destruí-lo mata todos os crânios vinculados.}
\end{figure}

\begin{description}
    \item[Lore:] Âncora do sistema imunológico de Érebo. O Totem absorve a matéria orgânica dos Goblins mortos nas minas subterrâneas e reanima seus crânios como projéteis rastreadores. Os Goblins não sabem — ou fingem não saber — que seus mortos servem ao planeta desta forma.
    \item[Comportamento:]
    \begin{itemize}
        \item Estático. Ativa quando o jogador entra em \texttt{20m}.
        \item Spawna até \textbf{3 HomingHazards} (crânios de Goblin reanimados) a cada \texttt{5s} em posições aleatórias dentro de \texttt{10m}.
        \item Os crânios spawnaados \textbf{não consomem budget} — são gerados pela habilidade do Totem, não pelo sistema de ondas.
        \item \textbf{Morte em Cadeia:} Ao ser destruído, todos os crânios vinculados são destruídos instantaneamente.
    \end{itemize}
    \item[Stats:] Crânios por ciclo \texttt{3} --- Intervalo \texttt{5s} (buffado \texttt{2.5s}) --- Raio \texttt{10m}
    \item[Buff do Sintonizador:] Intervalo de spawn $\div$2.
\end{description}

% =========================================================================
\section{Tabela Resumo --- Bioma 1 (Cristalizado)}

\begin{center}
\begin{tabular}{|l|l|l|}
    \hline
    \textbf{Inimigo} & \textbf{Classe de Spawn} & \textbf{Script} \\
    \hline
    Aranha               & Mob Menor (1 pt)          & \texttt{Spider\_AI} \\
    Goblin               & Atirador (4 pts)          & \texttt{GoblinAI\_Transform} \\
    Golem                & Tanque (4 pts)            & \texttt{Golem\_AI} \\
    Sintonizador         & Suporte (5 pts)           & \texttt{CrystalTuner} \\
    Totem                & Estrutura (6 pts)         & \texttt{TotemSpawner} \\
    Pedra Mágica         & Elite (10 pts)            & \texttt{MagicStone\_AI} \\
    Enxame de Cristais   & Elite (10 pts)            & \texttt{ShardSwarm\_AI} \\
    Vigia de Cristal     & Elite (10 pts)            & \texttt{CrystalWatcher\_AI} \\
    \hline
\end{tabular}
\end{center}

\begin{itemize}
    \item \textbf{Chefe do Bioma 1:} (a ser definido).
\end{itemize}

% =========================================================================
% CAPÍTULO 7: BIOMAS E EXPLORAÇÃO
% =========================================================================
\chapter{Biomas e Exploração}

O planeta é dividido em zonas distintas, mas o jogador só conhece a área da queda inicial. As outras áreas são desconhecidas e hostis.

\section{1. Bioma Cristalizado (Zona da Queda)}
\begin{itemize}
    \item \textbf{Visual:} Vastos campos abertos, cores frias (ciano, roxo), formações geológicas afiadas.
    \item \textbf{Inimigos:} Formas de vida baseadas em silício e cristal. Rápidos e cortantes. (Ver Capítulo~\ref{cap:bestiario} para detalhes de cada classe.)
    \item \textbf{Chefe:} (?)
\end{itemize}

\section{2. [DADOS CORROMPIDOS]}
\textit{Setor inexplorado. Análises preliminares indicam alta toxicidade atmosférica. A ser trabalhado futuramente.}

\section{3. [DADOS CORROMPIDOS]}
\textit{Setor inexplorado. Assinaturas energéticas de metal e carne detectadas. A ser trabalhado futuramente.}

\section{4. [ACESSO NEGADO PELA ENTIDADE]}
\textit{O Núcleo do Planeta. Leitura impossível.}

% =========================================================================
% CAPÍTULO 8: A NAVE DESTRUÍDA (BASE)
% =========================================================================
\chapter{A Nave Destruída (Base)}

Entre as runs, o jogador retorna aos destroços da \textit{Nave}. É o único local onde o solo foi ``esterilizado'' pelo combustível da nave, impedindo que o planeta a consuma.

\begin{itemize}
    \item \textbf{Bancada de Reparos:} Onde o jogador gasta os minerais que trouxe na bolsa para melhorar seu traje (Vida/Armadura) e a eficiência de suas ferramentas.
    \item \textbf{O Canil:} Onde Eptinho vive. Interagir com ele pode dar dicas sobre o próximo bioma ou ``Buffs de Moral'' (Sorte).
    \item \textbf{Objetivo Final:} Consertar os propulsores principais para escapar da gravidade do planeta.
\end{itemize}

% =========================================================================
% CAPÍTULO 9: GERAÇÃO PROCEDURAL DE MAPAS
% =========================================================================
\chapter{Geração Procedural de Mapas}\label{cap:geracao}

O mapa de cada \textit{run} é gerado dinamicamente no início da partida. O sistema é baseado em encaixe de \textbf{Prefabs de Salas} via um algoritmo de \textbf{Sockets Direcionais}.

\section{Visão Geral da Arquitetura}
\begin{description}
    \item[\texttt{GameManager}:] Gerencia o ciclo de vida das cenas. Ao chamar \texttt{LoadGameLevel()}, dispara o carregamento assíncrono e aciona o \texttt{LevelGenerator}.
    \item[\texttt{LevelGenerator}:] Responsável pela montagem procedural do mapa via sockets. Notifica o \texttt{GameManager} quando o nível está pronto.
    \item[\texttt{RoomController}:] Presente em cada sala de combate. Gerencia ondas de inimigos e trava/destrava portas.
    \item[\texttt{RunManager}:] Persistente entre cenas (\textit{DontDestroyOnLoad}). Rastreia o índice de sala e fornece parâmetros econômicos de dificuldade.
\end{description}

\section{O Sistema de Sockets}
\begin{center}
\begin{tabular}{|c|c|c|}
    \hline
    \textbf{Nome do Socket} & \textbf{Direção} & \textbf{Uso} \\
    \hline
    \texttt{North} & \texttt{Vector3.forward} & Saída para o Norte \\
    \texttt{South} & \texttt{Vector3.back}    & Saída para o Sul  \\
    \texttt{East}  & \texttt{Vector3.right}   & Saída para o Leste \\
    \texttt{West}  & \texttt{Vector3.left}    & Saída para o Oeste \\
    \texttt{Entrance} & Variável & Entrada de corredores e salas especiais \\
    \hline
\end{tabular}
\end{center}

\section{Tipos de Sala e Prefabs}
\begin{center}
\begin{tabular}{|l|c|l|}
    \hline
    \textbf{Prefab} & \textbf{Dimensões} & \textbf{Função} \\
    \hline
    \texttt{Room\_20x20\_Safe}       & 20×20 & Sala de início (segura, sem inimigos) \\
    \texttt{Room\_20x40}             & 20×40 & Sala principal de combate (média) \\
    \texttt{Room\_20x60}             & 20×60 & Sala principal de combate (grande) \\
    \texttt{Room\_10x20\_Transition} & 10×20 & Corredor de transição entre salas \\
    \texttt{Room\_Merchant}          & —     & Beco sem saída — possível presença do Mercador \\
    \texttt{ExitRoom\_Prefab}        & —     & Sala com portão para o próximo nível \\
    \texttt{Room\_DeadEnd}           & —     & Beco sem saída selado (sem evento) \\
    \hline
\end{tabular}
\end{center}

\section{Regras de Geração}\label{sec:regras_geracao}
\begin{description}
    \item[Quantidade máxima de salas (\texttt{maxMainRooms}):] O gerador tenta construir até \textbf{10 salas de combate} por \textit{run}.
    \item[Limite por tipo (\texttt{roomLimitPerType}):] Cada prefab pode aparecer no máximo \textbf{5 vezes} por \textit{run}.
    \item[Sala Inicial:] A \texttt{Room\_20x20\_Safe} é sempre a primeira. O algoritmo só abre sockets que não sejam \texttt{West} a partir dela.
    \item[Sala de Saída:] Após o loop principal, \textbf{um} dos becos sem saída é aleatoriamente escolhido para receber a \texttt{ExitRoom}.
    \item[Sala do Mercador:] Cada beco sem saída restante tem \textbf{25\% de chance} de receber o Mercador. \textbf{Não há garantia de fallback}: se nenhum beco sortear sucesso, o Mercador simplesmente não aparece nessa \textit{run}. Encontrá-lo é um evento genuinamente probabilístico.
    \item[Becos Restantes:] Todos os demais becos são selados com a \texttt{Room\_DeadEnd}.
\end{description}

\section{Fluxo de Geração (Resumo)}
\begin{enumerate}
    \item \textbf{Sala Inicial} é instanciada em \texttt{(0,0,0)}. Sockets abertos são adicionados à \textit{frontier}.
    \item \textbf{Loop Principal:} Enquanto há salas abaixo de \texttt{maxMainRooms} e sockets disponíveis:
    \begin{enumerate}
        \item Consome um socket aleatório da \textit{frontier}.
        \item Instancia corredor de transição e encaixa.
        \item Escolhe sala principal válida e encaixa.
        \item Adiciona os sockets da nova sala à \textit{frontier}.
    \end{enumerate}
    \item \textbf{Pós-Loop:} Distribui sala de saída, tenta spawnar Mercador (25\% por beco), sela becos restantes.
    \item \textbf{Spawn de Itens:} O \texttt{ItemSpawner} é acionado para popular o mapa.
    \item \textbf{Jogador é Posicionado:} O \texttt{GameManager} posiciona o jogador no \texttt{Player\_StartPoint} da sala inicial.
\end{enumerate}

% =========================================================================
% CAPÍTULO 10: SISTEMA ECONÔMICO (IMPLEMENTAÇÃO)
% =========================================================================
\chapter{Sistema Econômico (Implementação)}\label{cap:economia}

Este capítulo detalha as fórmulas e parâmetros do sistema econômico \textbf{atualmente implementados} no código, servindo como referência técnica para a equipe de design de níveis e balanceamento.

\section{Progressão de Dificuldade das Salas}

\subsection{Fórmula de Budget de Spawn}
\[
P(n) = B_{spawn} + G_{spawn} \times n
\]
Onde $B_{spawn} = 10$, $G_{spawn} = 0{,}9$, e $n$ = índice da sala (começa em 1).

\begin{center}
\begin{tabular}{|c|c|l|}
    \hline
    \textbf{Sala} & \textbf{Budget Total} & \textbf{Composição possível (por onda)} \\
    \hline
    1  & 11 pts & $\approx$ 11 Mob Menores \\
    5  & 15 pts & 1 Elite (10) + 5 Mob Menores \\
    10 & 19 pts & 1 Elite + 2 Tanques + 1 Mob Menor \\
    20 & 28 pts & 1 Elite + 4 Tanques + 2 Atiradores + extras \\
    \hline
\end{tabular}
\end{center}

\subsection{Budget por Onda}
\[
\text{Budget}_{onda} = \max\!\left(4,\ \left\lfloor \frac{P(n)}{\text{Total de ondas}} \right\rceil \right)
\]

\subsection{Fórmula de Drop de Essência}
\[
E(n) = d \times (1 + \alpha_{drop} \times n)
\]
Onde $d$ é o drop base do inimigo e $\alpha_{drop} = 0{,}05$ (+5\% por sala).

\section{Classes de Inimigos e Sistema de Pontos}
\begin{center}
\begin{tabular}{|l|c|l|}
    \hline
    \textbf{Classe} & \textbf{Custo} & \textbf{Restrições} \\
    \hline
    Mob Menor              & 1 ponto   & Máx. de 50--70\% dos pontos da onda \\
    Atirador               & 4 pontos  & Sempre em par (2 unidades por 4 pts) \\
    Tanque                 & 4 pontos  & Máx. de 4 por sala (limite global) \\
    Suporte (Sintonizador) & 5 pontos  & Máx. de 2 por sala \\
    Estrutura (Totem)      & 6 pontos  & Máx. de 1 por sala; crânios não consomem budget \\
    Elite                  & 10 pontos & Máx. de 1 por sala (limite global) \\
    \hline
\end{tabular}
\end{center}

\subsection{Garantia de Mínimo de Inimigos}
Independentemente do budget, o \texttt{RoomController} garante ao menos \texttt{minEnemiesPerWave = 6} inimigos por onda, preenchendo com \textbf{Mob Menores} sem custo adicional.

\section{Sistema de Ondas por Sala}
\begin{itemize}
    \item \textbf{Número de ondas:} Sorteado entre \texttt{minWaves = 1} e \texttt{maxWaves = 3}.
    \item \textbf{Intervalo entre ondas:} \texttt{timeBetweenWaves = 2} segundos.
    \item \textbf{Delay pré-onda:} 1 segundo após as portas travarem.
    \item \textbf{Indicador Visual:} Todo inimigo é gerado com um marcador visual (\texttt{SpawnIndicator}) que aparece 1,5 segundo antes do inimigo surgir.
\end{itemize}

\section{Sistema de Infusão e Inflação de Custo}\label{sec:inflacao}

A fórmula de custo de infusão implementada no \texttt{InfusionManager} é:
\[
C = B \times (1{,}0 + \alpha \times P_{total})
\]
Onde:
\begin{itemize}
    \item $B$ = custo base do item (definido pelo Tier)
    \item $\alpha = 0{,}1$ \quad (coeficiente de inflação por \textit{run})
    \item $P_{total}$ = soma dos \textbf{pesos de inflação} de todos os itens \textbf{infundidos} nesta \textit{run}
\end{itemize}

\subsection{Custos Base e Pesos por Tier}
\begin{center}
\begin{tabular}{|c|l|c|c|c|}
    \hline
    \textbf{Tier} & \textbf{Nome} & \textbf{Cor} & \textbf{Custo Base ($B$)} & \textbf{Peso ($P$)} \\
    \hline
    T1 & Comum    & Branco/Cinza & 60  & 1,00 \\
    T2 & Incomum  & Verde        & 180 & 2,25 \\
    T3 & Raro     & Azul         & 300 & 4,00 \\
    T4 & Lendário & Dourado      & 420 & 6,00 \\
    \hline
\end{tabular}
\end{center}

\subsection{Regras Críticas do $P_{total}$}

\begin{description}
    \item[Infusão Normal:] Ao infundir um item (de qualquer origem), seu peso $P$ é somado ao $P_{total}$.
    \item[Cirurgia de Remoção (Mercador):] Ao remover cirurgicamente um item já infundido, o peso $P$ daquele item é \textbf{subtraído} do $P_{total}$. O $P_{total}$ nunca vai abaixo de zero. Exemplo: se $P_{total} = 5{,}25$ e um T2 ($P = 2{,}25$) é removido, o novo $P_{total} = 3{,}00$.
    \item[Compra no Mercador (Artefatos):] Comprar um item identificado no Mercador \textbf{não altera} o $P_{total}$. O $P_{total}$ só aumenta no momento da \textbf{infusão} do item, independente de sua origem.
\end{description}

\subsection{Exemplo de Escalonamento de Custo}
Considere um jogador que infundiu sequencialmente: T1, T1, T2.
\[
P_{total} = 1{,}00 + 1{,}00 + 2{,}25 = 4{,}25
\]
O custo da próxima infusão de um item T3 ($B = 300$):
\[
C = 300 \times (1{,}0 + 0{,}1 \times 4{,}25) = 300 \times 1{,}425 = \textbf{428 Essências}
\]
Um aumento de \textbf{+42,5\%} sobre o custo original.

\section{Criação de Itens (Pipeline de Design)}
Cada item do jogo é um \textbf{ScriptableObject} do tipo \texttt{ItemData}, criado via \texttt{Assets > Create > Inventory > Item Data}.

\begin{description}
    \item[\texttt{itemId}:] Identificador único em texto. Deve casar com o \texttt{CharacteristicItemPickup} no prefab inimigo.
    \item[\texttt{tier}:] O \textit{Tier} de raridade (Common / Uncommon / Rare / Legendary).
    \item[\texttt{infusionEssenceCost}:] Custo base $B$ (recomendado: 60 / 180 / 300 / 420 por tier).
    \item[\texttt{itemAttributes}:] Lista de \texttt{ItemAttributeParam}, onde cada entrada define \textbf{qual atributo} é modificado, \textbf{quanto}, e se o valor é \textbf{somado} ou \textbf{multiplicado}.
\end{description}

\subsection{Atributos Disponíveis para Itens}
\begin{center}
\begin{tabular}{|l|l|}
    \hline
    \textbf{Categoria} & \textbf{Atributos} \\
    \hline
    Ofensivos  & \texttt{BaseDamageMultiplier}, \texttt{AttackSpeedMelee}, \texttt{CritChance}, \\
               & \texttt{CritMultiplier}, \texttt{Knockback}, \texttt{WeaponRangeMelee}, \\
               & \texttt{WeaponRangeProjectile}, \texttt{Piercing}, \texttt{BounceChance}, \\
               & \texttt{BounceCount}, \texttt{MultiShotChance}, \texttt{Spread} \\
    \hline
    Defensivos & \texttt{MaxHealth}, \texttt{MaxArmor}, \texttt{ArmorRegen}, \\
               & \texttt{DodgeChance}, \texttt{DamageNegation}, \texttt{Thorns} \\
    \hline
    Mobilidade & \texttt{SpeedMultiplier}, \texttt{DashCooldownMultiplier}, \\
               & \texttt{DashCounts}, \texttt{DashInvulnerability} \\
    \hline
\end{tabular}
\end{center}

\end{document}
